\section{Neural ODE}

The neural ODE \cite{chen2019neuralordinarydifferentialequations} substitutes the discrete neural network hidden layers by a differentiable function. 
It is a smooth version of residual neural networks. 
The output of this continuous-depth network is the solution of an initial value point ODE problem ending at an arbitrarily depth. 
The forward propagation does not store the gradients and intermediate values and is conducted by a black-box ODE solver. 
The neural ODE based models have constant memory cost and can trade numerical precision with time cost. 

\subsection{Method}

The continuous-depth network is described by an initial value point problem with an invariant network function $f$, as shown below. 
We suppose the output of latent state $z$ and network function $f$ are all vectors. 
\begin{equation*}
    \left\{\begin{aligned}
        &\frac{dz(t)}{dt} = f(z(t),t,\theta) \\
        &z(t_0) = z_0
    \end{aligned}\right.
\end{equation*}

The output at time $t_1$ is numerically solved by a black-box ODE solver, for example Runge-Kutta method. 
Analytically, we have the solution given by 
\begin{equation*}
    z(t_1) = z(t_0) + \int_{t_0}^{t_1} f(z(t),t,\theta)dt. 
\end{equation*}
The corresponding loss of the model at time $t_1$ is 
\begin{equation*}
    L(z(t_1)) = L\left(z(t_0) + \int_{t_0}^{t_1} f(z(t),t,\theta)dt\right). 
\end{equation*}
The gradient w.r.t the model parameter $\theta$ is calculated by 
\begin{align*}
    \frac{\partial L}{\partial \theta} = \left(\frac{\partial z(t_1)}{\partial \theta}\right)^T\frac{\partial L}{\partial z(t_1)} &= \left(\frac{\partial z(t_0)}{\partial \theta}+\frac{\partial}{\partial \theta}\int_{t_0}^{t_1}f(z(t),t,\theta)dt\right)^T\frac{\partial L}{\partial z(t_1)} \\
                                                                                                                   &= \left(\frac{\partial z(t_0)}{\partial \theta}+\int_{t_0}^{t_1}\frac{\partial f}{\partial \theta}+\frac{\partial f}{\partial z(t)}\frac{\partial z(t)}{\partial \theta}dt\right)^T\frac{\partial L}{\partial z(t_1)}, 
\end{align*}
where the gradients $\frac{\partial L}{\partial z(t)}$ and $\frac{\partial f}{\partial z(t)}$ can be calculated by auto-differentiation since $f$ is clearly known.
However, $\frac{\partial z(t)}{\partial \theta}$ is intractable since the latent dynamic $z(t)$ is unobservable.

To solve this, we introduce the adjoint sensitivity method to get the trajectory of $\frac{\partial L}{\partial z(t)}$ by solving another ODE. 
This helps removing the problematic gradients in computing $\frac{\partial L}{\partial \theta}$. 

\subsubsection{The adjoint and two equations}

We define the adjoint by $a(t):=\frac{\partial L}{\partial z(t)}$. 
The adjoint satisfies two equations 
\begin{align}
    \frac{da(t)}{dt} &= -\left(\frac{\partial}{\partial z(t)}f(z(t),t,\theta)\right)^Ta(t), \label{eq: sec2_adjoint1} \\
    \frac{\partial L}{\partial \theta} &= -\int_{t_1}^{t_0}\left(\frac{\partial f}{\partial \theta}\right)^Ta(t)dt,  \label{eq: sec2_adjoint2}
\end{align}
where the gradient $\frac{\partial f}{\partial \theta}$ in (\ref{eq: sec2_adjoint2}) is the derivative of $f(z(t),t,\theta)$ w.r.t $\theta$ treating $z$ as constant. 
The two equations record the process of calculating the prior knowledge for updating the parameter. 
(\ref{eq: sec2_adjoint2}) depends on the trajectory of $a(t)$ which means we need to simultaneously compute the value of $a(t)$ when calculating the numerical approximation of the integral. 
Before understanding the full backward propagation process, 
we give the proofs of the two equations respectively. 

\paragraph{Proof of (\ref{eq: sec2_adjoint1})}
The proof relies on using the definition of derivative as $t_1$ and $t_0$ are arbitrarily close. 
We start by computing the derivative of loss about $z(t_0)$. 

Suppose $a(t):=\frac{\partial L}{\partial z(t)}$ is differentiable about $t$. 
We consider a scalar-valued loss $L$. 
\begin{align*}
    \frac{\partial L}{\partial z(t_0)} &= (\frac{\partial z(t_1)}{\partial z(t_0)})^T\frac{\partial L}{\partial z(t_1)} \\
                                       &= \left(I+\frac{\partial}{\partial z(t_0)}\int_{t_0}^{t_1} f(z(t),t,\theta)dt\right)^T\frac{\partial L}{\partial z(t_1)} \\
                                       &= \frac{\partial L}{\partial z(t_1)} + \left(\frac{\partial}{\partial z(t_0)}\int_{t_0}^{t_1} f(z(t),t,\theta)dt\right)^T\frac{\partial L}{\partial z(t_1)} \\
    \frac{\partial L}{\partial z(t_1)}-\frac{\partial L}{\partial z(t_0)} &= -\left(\frac{\partial}{\partial z(t_0)}\int_{t_0}^{t_1} f(z(t),t,\theta)dt\right)^T\frac{\partial L}{\partial z(t_1)} \\
    \frac{\frac{\partial L}{\partial z(t_1)}-\frac{\partial L}{\partial z(t_0)}}{t_1-t_0} &= -\frac{1}{t_1-t_0}\left(\frac{\partial}{\partial z(t_0)}\int_{t_0}^{t_1} f(z(t),t,\theta)dt\right)^T\frac{\partial L}{\partial z(t_1)} \\
    \lim_{t_1\rightarrow t_0} \frac{\frac{\partial L}{\partial z(t_1)}-\frac{\partial L}{\partial z(t_0)}}{t_1-t_0} &= -\lim_{t_1\rightarrow t_0}\frac{1}{t_1-t_0}\left(\frac{\partial}{\partial z(t_0)}\int_{t_0}^{t_1} f(z(t),t,\theta)dt\right)^T\frac{\partial L}{\partial z(t_1)} \\
    \frac{\partial}{\partial t}\frac{\partial L}{\partial z(t)}\bigg|_{t=t_0} &= -\left(\frac{\partial}{\partial z(t_0)}f(z(t_0),t,\theta)\right)^T\frac{\partial L}{\partial z(t_0)} \\
    \frac{da(t)}{d(t)}\bigg|_{t=t_0} &= -\left(\frac{\partial}{\partial z(t_0)}f(z(t_0),t,\theta)\right)^Ta(t_0).
\end{align*}
By the randomness of $t_0$, we have 
\begin{equation*}
    \frac{da(t)}{dt} = -\left(\frac{\partial}{\partial z(t)}f(z(t),t,\theta)\right)^Ta(t).
\end{equation*}

\paragraph{Proof of (\ref{eq: sec2_adjoint2})}
Aside from the definition of derivative, the proof relies on the trajectory of adjoint given by the Equation \ref{eq: sec2_adjoint1}. 
We also need the intermediate expression of the gradient $\frac{\partial L}{\partial\theta}$. 
We suppose the parameter $\theta\in\mathbb{R}^d$.

\begin{align*}
    \frac{\partial L}{\partial \theta} = \left(\frac{\partial z(t_1)}{\partial \theta}\right)^Ta(t_1) &= \left(\frac{\partial z(t_0)}{\partial \theta}+\frac{\partial}{\partial \theta}\int_{t_0}^{t_1}f(z(t),t,\theta)dt\right)^Ta(t_1) \\
                                                                                                                   &= \left(\frac{\partial z(t_0)}{\partial \theta}+\int_{t_0}^{t_1}\frac{\partial f}{\partial \theta}+\frac{\partial f}{\partial z(t)}\frac{\partial z(t)}{\partial \theta}dt\right)^Ta(t_1). 
\end{align*}
From the above deduction, we have a new equation.
\begin{align*}
    \left(\frac{\partial z(t_1)}{\partial \theta}\right)^Ta(t_1) &= \left(\frac{\partial z(t_0)}{\partial \theta}+\int_{t_0}^{t_1}\frac{\partial f}{\partial \theta}+\frac{\partial f}{\partial z(t)}\frac{\partial z(t)}{\partial \theta}dt\right)^Ta(t_1) \\
    (\frac{\partial z(t_1)}{\partial \theta} - \frac{\partial z(t_0)}{\partial \theta})^Ta(t_1) &= \left(\int_{t_0}^{t_1}\frac{\partial f}{\partial \theta}+\frac{\partial f}{\partial z(t)}\frac{\partial z(t)}{\partial \theta}dt\right)^Ta(t_1) \\
    \lim_{t_1\rightarrow t_0}\frac{1}{t_1-t_0}(\frac{\partial z(t_1)}{\partial \theta} - \frac{\partial z(t_0)}{\partial \theta})^Ta(t_1) &= \lim_{t_1\rightarrow t_0}\frac{1}{t_1-t_0}\left(\int_{t_0}^{t_1}\frac{\partial f}{\partial \theta}+\frac{\partial f}{\partial z(t)}\frac{\partial z(t)}{\partial \theta}dt\right)^Ta(t_1) \\
    \left(\frac{\partial}{\partial t}\frac{\partial z(t)}{\partial\theta}\right)^T\bigg|_{t=t_0}a(t_0) &= \left(\frac{\partial f}{\partial \theta}+\frac{\partial f}{\partial z(t)}\frac{\partial z(t)}{\partial \theta}\right)^T\bigg|_{t=t_0}a(t_0).
\end{align*}
By the randomness of $t_0$, we have 
\begin{align*}
    \left(\frac{\partial}{\partial t}\frac{\partial z(t)}{\partial\theta}\right)^Ta(t) = \left(\frac{\partial f}{\partial \theta}+\frac{\partial f}{\partial z(t)}\frac{\partial z(t)}{\partial \theta}\right)^Ta(t) &= \left(\frac{\partial f}{\partial \theta}\right)^Ta(t)+ \left(\frac{\partial f}{\partial z(t)}\frac{\partial z(t)}{\partial \theta}\right)^Ta(t) \\
                                                                                       &= \left(\frac{\partial f}{\partial \theta}\right)^Ta(t)+\left(\frac{\partial z(t)}{\partial \theta}\right)^T\left(\frac{\partial f}{\partial z(t)}\right)^Ta(t) \\
                                                                                       &= \left(\frac{\partial f}{\partial \theta}\right)^Ta(t) - \left(\frac{\partial z(t)}{\partial \theta}\right)^T\frac{da(t)}{dt} & (\text{By (\ref{eq: sec2_adjoint1})}) \\
    \left(\frac{\partial}{\partial t}\frac{\partial z(t)}{\partial\theta}\right)^Ta(t) + \left(\frac{\partial z(t)}{\partial \theta}\right)^T\frac{da(t)}{dt} &= \left(\frac{\partial f}{\partial \theta}\right)^Ta(t) \\
    \frac{d}{dt}\left(\left(\frac{\partial z(t)}{\partial\theta}\right)^Ta(t)\right) &= \left(\frac{\partial f}{\partial \theta}\right)^Ta(t) \\
    \left(\frac{\partial z(t_1)}{\partial\theta}\right)^Ta(t_1) - \left(\frac{\partial z(t_0)}{\partial\theta}\right)^Ta(t_0) &= \int_{t_0}^{t_1}\left(\frac{\partial f}{\partial \theta}\right)^Ta(t)dt \\
    \left(\frac{\partial z(t_1)}{\partial\theta}\right)^Ta(t_1) &= \int_{t_0}^{t_1}\left(\frac{\partial f}{\partial \theta}\right)^Ta(t)dt. 
\end{align*}
The last step is by the fact that the initial layer $z(t_0)$ is irrelevant to the parameter $\theta$.  
Since the LHS is just the gradient $\frac{\partial L}{\partial \theta}$, 
we have 
\begin{equation*}
    \frac{\partial L}{\partial \theta} = -\int_{t_1}^{t_0}\left(\frac{\partial f}{\partial \theta}\right)^Ta(t)dt. 
\end{equation*}

\subsubsection{Backward propagation of neural ODE}

To get the gradient $\frac{\partial L}{\partial\theta}$, 
we compute the integral (\ref{eq: sec2_adjoint2}) and need the simultaneous computation of the adjoint $a(t)$ through (\ref{eq: sec2_adjoint1}). 
Among the two processes, the gradients $\frac{\partial f}{\partial z(t)}$ and $\frac{\partial f}{\partial\theta}$ are calculated by auto-differentiation 
and the vector-Hessian products are computed subsequently. 
The adjoint trajectory ODE (\ref{eq: sec2_adjoint1}) is solved by a random ODE-solver starting from the initial time $t_1$ and initial value $a(t_1)=\frac{\partial L}{\partial z(t_1)}$. 
At each end of a step, the output $a(t_{intermediate})$ is stored temporarily and accumulated by the integral. 
This proceeds until $t_0$, making sure the memory cost is a constant. 